% Part: sets-functions-relations
% Chapter: functions
% Section: functions-relations

\documentclass[../../../include/open-logic-section]{subfiles}


\begin{document}

\olfileid{sfr}{fun}{rel}

\olsection{فنکشن تعلقاں دے طور تے}

\begin{explain}
اوہ فنکشن جیہڑا~$A$ دے !!{element}s نوں~$B$ دے !!{element}s نال
جوڑدا اے، ظاہر اے کہ $A$ تے~$B$ دے وچکار اک تعلق تعریف کردا اے:
اوہی تعلق جیہڑا $x$ تے $y$ دے وچکار عین اوس ویلے قائم اے جدوں
$f(x) = y$ ہووے۔ اصل وچ، جے مثال لئی ساڈی دلچسپی ریاضی دے بنیادی
اجزا گھٹاؤن وچ ہووے، تاں اسی فنکشن~$f$ نوں ایس تعلق نال، یعنی جوڑیاں
دے اک سیٹ نال، \emph{اکو سمجھ} وی سکدے آں۔ فیر ایہ سوال پیدا ہوندا
اے: کیہڑے تعلق ایس طریقے نال فنکشن تعریف کردے نیں؟
\end{explain}

\begin{defn}[فنکشن دا گراف] من لو کہ $f\colon A \to B$ اک فنکشن اے۔
فنکشن~$f$ دا \emph{گراف} اوہ تعلق $R_f \subseteq A \times B$ اے
جیہڑا ایویں تعریف ہوندا اے:
\[
R_f = \Setabs{\tuple{x,y}}{f(x) = y}.
\]
\end{defn}

\begin{explain}
عنصراں نال تعیّن دے اصول کرکے فنکشن دا گراف اکلوتے طور تے طے ہوندا
اے۔ ہور ایہ کہ، سیٹاں اُتے عنصراں نال تعیّن دا اصول فنکشناں لئی
ضمنی اصول نوں فوراً ثابت کردا اے: جے $f$ تے~$g$ دا دائرۂ تعریف تے
ہدف سیٹ سانجھا ہووے، تاں ہر قدر اُتے متفق ہون دی صورت وچ اوہ اکو
فنکشن نیں۔

ایسے طرح، جے کوئی تعلق ``فنکشنی'' ہووے، تاں اوہ کسے فنکشن دا گراف اے۔
\end{explain}

\begin{prop}\ollabel{prop:graph-function}
من لو کہ $R \subseteq A \times B$ ایس طرح اے کہ:
\begin{enumerate}
\item جے $Rxy$ تے $Rxz$، تاں $y = z$؛ تے
\item ہر $x \in A$ لئی کوئی $y \in B$ اے جیہدے لئی $\tuple{x,
y} \in R$۔
\end{enumerate}
فیر $R$ اوس فنکشن $f\colon A \to B$ دا گراف اے جیہڑا
$f(x) = y$ عین اوس ویلے جدوں $Rxy$ راہیں تعریف ہوندا اے۔
\end{prop}

\begin{proof}
من لو کہ کوئی $y$ اے جیہدے لئی $Rxy$۔ جے کوئی ہور $z \neq y$ وی
ہوندا جیہدے لئی $Rxz$، تاں~$R$ اُتے شرط دی خلاف ورزی ہوندی۔ ایس
لئی جے کوئی $y$ اے جیہدے لئی $Rxy$، تاں ایہ $y$ اکلوتا اے، تے ایس
واسطے $f$ صحیح طور تے تعریف شدہ اے۔ ظاہر اے کہ $R_f = R$۔
\end{proof}

\begin{explain}
ہر فنکشن $f\colon A \to B$ دا اک گراف ہوندا اے، یعنی دائرۂ تعریف
تے ہدف سیٹ دے وچکار اک تعلق (برابر طور تے $A \times B$ دا اک ذیلی سیٹ) جیہڑا
$f(x) = y$ راہیں تعریف ہوندا اے۔ دوجے پاسے، ہر تعلق~$R
\subseteq A \times B$ جیہدے وچ \olref{prop:graph-function} والیاں
خاصیتاں ہون، اک فنکشن~$f \colon A \to B$ دا گراف اے۔ فنکشناں تے
اوہناں دے گرافاں دے ایس نیڑے رشتے کرکے اسی فنکشن نوں بس اوہدا گراف
سمجھ سکدے آں۔ ہور لفظاں وچ، فنکشناں نوں کجھ مخصوص تعلقاں، یعنی
ٹپلاں دے کجھ مخصوص سیٹاں نال شناخت کیتا جا سکدا اے۔
% Source correction OLFUN-004: relation between A and B / subset of A times B, not a relation on A times B.
\oliflabeldef{sfr:rel:ref:sec}{پر ایہ یاد رکھو کہ ایس ``شناخت'' دی
روح \olref[sfr][rel][ref]{sec} والی اے: ایہ فنکشناں دی مابعد الطبیعی
حقیقت بارے دعویٰ نہیں، بلکہ ایہ مشاہدہ اے کہ فنکشناں نوں کجھ سیٹاں
دے طور تے \emph{ورتنا} سوکھا اے۔ ایس سہولت دی اک وجہ ایہ اے کہ }{}
ہُن اسی فنکشناں اُتے وی اوہو جہے عمل کر سکدے آں جیہڑے اسی تعلقاں
اُتے کیتے سن (ویکھو \olref[sfr][rel][ops]{sec})۔ خاص طور تے:
\end{explain}

\begin{defn}\ollabel{defn:funimage}
من لو کہ $f \colon A \to B$ اک فنکشن اے تے $C\subseteq A$۔

فنکشن~$f$ دی~$C$ تک \emph{تحدید} اوہ
فنکشن~$\funrestrictionto{f}{C}\colon C \to B$ اے جیہڑا
$(\funrestrictionto{f}{C})(x) = f(x)$ راہیں ہر $x \in C$ لئی
تعریف ہوندا اے۔ ہور لفظاں وچ، $\funrestrictionto{f}{C} =
\Setabs{\tuple{x, y} \in R_f}{x \in C}$۔

فنکشن~$f$ دا~$C$ اُتے \emph{اطلاق} $\funimage{f}{C} =
\Setabs{f(x)}{x \in C}$ اے۔ اسی ایہنوں~$C$ دا~$f$ دے تحت
\emph{عکس} وی آکھدے آں۔
\end{defn}

\begin{explain}
ایہناں تعریفاں توں نکلدا اے کہ $\ran{f} =
\funimage{f}{\dom{f}}$، ہر فنکشن~$f$ لئی۔
\oliflabeldef{sfr:rel:ops:sec}{ایہ خیال تعلقاں نال فنکشناں دی شناخت
تے \olref[sfr][rel][ops]{sec} دیاں تعریفاں نال اوہو طرح جُڑدے نیں
جِویں توقع اے۔ پر اتھے فنکشن دی تحدید صرف دائرۂ تعریف نوں ذیلی سیٹ
تک گھٹاؤندی اے تے ہدف سیٹ نوں اوہو رکھدی اے؛ ایس لئی ایہ پہلے
دتے یکساں سیٹ اُتے تعلق دی دو پاسیاں والی تحدید دا حرف بہ حرف
اوہی عمل نہیں۔ دوجے دو عمل---اُلٹے تے نسبتی حاصل ضرب---لئی کجھ ہور
تفصیل دی لوڑ اے۔ اسی اوہ \olref[inv]{sec} تے
\olref[cmp]{sec} وچ دیواں گے۔}{}
% Source correction OLFUN-005: domain-only function restriction; earlier homogeneous relation restriction is two-sided.
\end{explain}

\end{document}
